\section{Entitlement Mapping}\label{sec:entitlement-mapping}

\textit{Normative language per \href{https://tools.ietf.org/html/rfc2119}{RFC~2119}.}

\subsection{How Mapping Works --- Tenant Responsibility}

The tenant is solely responsible for the mapping between their users' entitlements
(expressed in tenant-internal identifiers) and Canvasflow resource identifiers.

Canvasflow resource identifiers are stable and self-served. The Canvasflow administration 
dashboard lists every publication and content resource the tenant is authorized to access 
alongside its \texttt{cf:<type>:<id>} identifier (for example, \texttt{cf:issue:123}). 
The tenant's technical lead logs into the dashboard, locates the relevant publications, and 
copies the identifiers directly from the UI. No mapping document is delivered by 
Canvasflow; no API is called to retrieve identifiers. This retrieval is a one-time 
step performed during Phase 2 of onboarding 
(see \hyperref[sec:canvasflow-side-setup]{Client On-Boarding \S9.2}).


From that point forward, maintaining and executing the mapping is entirely the tenant's
responsibility. The tenant maintains the mapping table inside their own IdP hook. The
form of the mapping table is the tenant's choice --- a hardcoded table, a group attribute
lookup, a call to the tenant's internal directory, or any other mechanism. The required
output format is defined in \hyperref[sec:mandatory-value-format]{\S4.2} below.

At every login \textbf{and} every token refresh, the hook reads the authenticated user's
entitlements from the tenant's source of truth, maps them to Canvasflow resource
identifiers, and embeds the result in the \textbf{access token} before it is signed.
Canvasflow never sees tenant-internal product identifiers and never performs mapping
logic.

\textbf{No Canvasflow endpoint is called at login time.} The mapping is performed
entirely inside the tenant's IdP hook without contacting Canvasflow's servers.

\subsection{Mandatory Value Format: the \texttt{cf:} Prefix}\label{sec:mandatory-value-format}

Every entitlement value emitted by the tenant's hook into the access token MUST use the
format:

\begin{lstlisting}
cf:<type>:<id>
\end{lstlisting}

where \texttt{<type>} is the Canvasflow resource type (e.g., \texttt{issue}) and
\texttt{<id>} is the stable Canvasflow-assigned numeric or alphanumeric identifier for
that resource. Examples of conforming values:

\begin{lstlisting}
cf:issue:1111
cf:issue:201
cf:issue:123
\end{lstlisting}

\subsubsection{The Ignore Rule for Non-Prefixed Values}

The Canvasflow Resource Server applies a strict filter to every value in the resolved
entitlement array:

\begin{itemize}
  \item Values with the \texttt{cf:} prefix and the well-formed
    \texttt{cf:<type>:<id>} structure --- \textbf{accepted} and used for content gating.
  \item \textbf{Any value without the \texttt{cf:} prefix --- silently ignored.} It is
    dropped from the effective entitlement set. It does not cause the token to be
    rejected. It does not grant any access.
\end{itemize}

The effective entitlement set is computed as:

\begin{lstlisting}
effective_entitlements =
  resolved_claim_array
    .filter(v => typeof v === 'string' &&
            v.startsWith('cf:'))
\end{lstlisting}

\textbf{Example:} a claim value of
\texttt{["cf:issue:123",\ "product:789",\ "issue:201"]} yields an effective entitlement
set of \texttt{["cf:issue:123"]}. The values \texttt{product:789} (a tenant-internal
identifier without the prefix) and \texttt{issue:201} (missing the \texttt{cf:} prefix)
are both ignored and logged server-side correlated by \texttt{jti}.

\textbf{Why the prefix is required.} The \texttt{cf:} namespace makes entitlement values
self-describing and collision-proof. A raw tenant-internal identifier such as \texttt{456}
or \texttt{product:789} that leaks into the claim can never be mistaken for a Canvasflow
entitlement. The prefix also enables unambiguous filtering --- the Resource Server requires
no knowledge of any tenant's internal identifier scheme to apply the rule.

\textbf{Why the ignore-rather-than-reject behavior.} IdP hooks and shared claims
frequently carry values from other integrations or from incomplete mapping configurations.
Rejecting an entire token because one value in the array lacks the \texttt{cf:} prefix
would turn a partial misconfiguration into a full authentication outage for the affected
user. Silently ignoring non-prefixed values fails safe: unmapped or foreign values grant
nothing, and the user still authenticates. Server-side logging of ignored values
(correlated by \texttt{jti}) preserves the ability for Canvasflow support to assist the
tenant in identifying and correcting the misconfiguration.

\subsection{Claim Name Hierarchy and Precedence}\label{sec:claim-name-hierarchy}

The mapped entitlement array MAY be communicated under any one of three claim names.
When more than one is present in a single token, the Resource Server applies a strict
precedence rule.




\begin{table}[htbp]
\centering
\begin{tabularx}{\textwidth}{@{}clX@{}}
\toprule
\textbf{Precedence} & \textbf{Claim name} & \textbf{Why it exists} \\
\midrule
1 (highest) & \texttt{cf:entitlements} &
  Namespaced, collision-proof. Some IdPs --- notably Auth0 --- blocklist non-namespaced
  custom claims on access tokens. \texttt{cf:entitlements} is guaranteed to work on
  every IdP and is the RECOMMENDED claim name for all new integrations. \\[6pt]
2 & \texttt{resources} &
  Legacy claim name retained for backward compatibility with tenants already integrated
  against this name. \\[6pt]
3 (lowest) & \texttt{entitlements} &
  The IANA-registered claim name per 
  \href{https://tools.ietf.org/html/rfc9068\#section-2.2.3.1}{RFC~9068~\S2.2.3.1} and 
  \href{https://tools.ietf.org/html/rfc7643\#section-4.1.2}{RFC~7643~\S4.1.2} ---
  standards-preferred but usable only on IdPs that do not reserve or blocklist it. \\
\bottomrule
\end{tabularx}
\caption{Entitlement claim name precedence}
\end{table}

\textbf{Resolution rule (normative):} The Resource Server reads the first claim present
in precedence order and uses it exclusively. It does not merge arrays across multiple
claim names. If \texttt{cf:entitlements} is present, \texttt{resources} and
\texttt{entitlements} are ignored entirely, even if they contain additional values.

\begin{lstlisting}
if token.claims["cf:entitlements"] exists:
    resolved_claim_array =
      token.claims["cf:entitlements"]
else if token.claims["resources"] exists:
    resolved_claim_array =
      token.claims["resources"]
else if token.claims["entitlements"] exists:
    resolved_claim_array =
      token.claims["entitlements"]
else:
    resolved_claim_array = []   
\end{lstlisting}

\textbf{Why precedence is inverted from standards purity.} The namespaced claim
\texttt{cf:entitlements} takes priority over the standards-registered
\texttt{entitlements} because \emph{reliability beats purity}. A tenant who deliberately
sets \texttt{cf:entitlements} for Canvasflow has expressed clear intent --- whereas the
bare \texttt{entitlements} claim may have been populated by the IdP itself (as a standard
SCIM attribute) or by an unrelated integration. When both are present, the namespaced
claim is the authoritative signal of what the tenant intended for this integration.

\textbf{Recommendation.} All new integrations SHOULD use \texttt{cf:entitlements}.
Existing integrations using \texttt{resources} or \texttt{entitlements} are supported
but SHOULD migrate to \texttt{cf:entitlements} at their next major hook update. Auth0
tenants MUST use \texttt{cf:entitlements} --- no other option is available (see
\hyperref[sec:per-idp-implementation]{Per-IdP Implementation Table \S6.2}).

\pagebreak

\subsection{Empty, Absent, and All-Ignored Claims: the Free-Content Baseline}\label{sec:free-content-baseline}

The following three states all resolve to the same effective condition:

\begin{table}[htbp]
\centering
\begin{tabularx}{\textwidth}{@{ }p{1.3cm}Xp{3cm}@{}}
\toprule
\textbf{State} & \textbf{Cause} & \textbf{Effective entitlement set} \\
\midrule
Empty array & \texttt{"cf:entitlements":[]} emitted by hook & \texttt{[]} \\[4pt]
Absent claim & None of the three claim names present & \texttt{[]} (logged as possible misconfiguration) \\[4pt]
All-ignored array & All values present but none carry \texttt{cf:} prefix & \texttt{[]} (logged as likely mapping misconfiguration) \\
\bottomrule
\end{tabularx}
\caption{Free-content baseline: three equivalent states}
\end{table}

In all three cases: \textbf{the user authenticates successfully and is entitled to free
content only.} Authentication does not fail. Content Canvasflow classifies as free is
served. Content Canvasflow classifies as premium is withheld. This is the semantic model
--- not an error condition.

Entitlements gate \emph{premium} content. Free content requires only a valid, verified
access token; entitlement values are not consulted.

\subsubsection{Consequences}

The following specific behaviors MUST be understood by tenant implementors:

\begin{itemize}
  \item \textbf{An empty array is a valid, intentional configuration.} A registered user
    on a free tier SHOULD have the hook emit \texttt{"cf:entitlements": []} explicitly.
    Explicit presence distinguishes ``free-tier user with correctly-configured hook'' from
    ``misconfigured hook with missing claim'' in server-side debugging.

  \item \textbf{An absent claim is tolerated} --- the behavior is identical to an empty
    array --- but Canvasflow MUST log a warning server-side (correlated by \texttt{jti})
    flagging the absence as a possible hook misconfiguration. An absent claim means
    Canvasflow cannot distinguish a deliberate free-tier configuration from a broken hook.

  \item \textbf{A present-but-non-array value MUST cause token rejection (401).} If the
    resolved claim name is present but its value is not a JSON array (e.g., it is a
    string, number, or object), the token is malformed and MUST be rejected. A malformed
    claim is a broken token, not a free-tier user.

  \item \textbf{An all-ignored array} (values present, none carry the \texttt{cf:}
    prefix) results in free content only. The non-prefixed values are logged server-side
    (correlated by \texttt{jti}) as a likely mapping misconfiguration, but the token is
    not rejected.
\end{itemize}

\pagebreak

\subsection{Why Mapping Happens at Auth Time in the Hook}

Mapping is executed in the IdP hook at authentication and refresh time rather than in
the Canvasflow GraphQL server at request time. The justification is as follows:

\textbf{Performance.} The hook is called once per login or refresh, not on every API
request. Executing entitlement mapping inside the GraphQL resolver would add a
cross-system lookup to every query, multiplying the latency cost proportionally with
API usage volume.

\textbf{Cryptographic integrity.} The result of the hook's mapping is sealed inside the
JWT's signature before the token leaves the tenant's IdP. Any post-issuance tampering
with the entitlement claim invalidates the signature. The Resource Server does not need
to trust any runtime claim about entitlements --- it trusts only the signed token.

\textbf{Resource Server purity.} The Canvasflow GraphQL server has no knowledge of any
tenant's internal product identifiers. The mapping result arrives as Canvasflow resource
identifiers (\texttt{cf:issue:123}) already resolved; the server applies a set membership
check --- not a mapping operation. This keeps the Resource Server stateless.

\textbf{Trade-off.} Entitlement changes take effect at the user's next login or refresh.
Existing tokens are valid until their \texttt{exp}. The maximum staleness window is
bounded by $\texttt{exp} \le 3600$ seconds. This trade-off is acceptable for
content delivery; if a tenant requires immediate revocation, the mitigation is to revoke
the user's refresh token in the IdP, which prevents new access tokens from being issued
and limits the exposure window to the remaining lifetime of the current access token.

\subsection{Why Tenant-Managed Rather Than a Canvasflow Mapping Endpoint}

An alternative design --- where Canvasflow provides an endpoint that the IdP hook calls
at login time to perform the mapping --- was considered and rejected:

\begin{itemize}
  \item It would add a runtime dependency on Canvasflow's availability to every user
    login. A Canvasflow outage would prevent tenants from logging in.
  \item It would require the IdP hook to hold a Canvasflow API key, creating a
    secret-management obligation on the tenant side.
  \item The mapping logic lives closest to the data it maps. The tenant knows their own
    product identifiers and their own users' entitlements; Canvasflow does not.
  \item The design would require Canvasflow to build, operate, and secure a
    tenant-facing mapping API --- increasing the attack surface of the Canvasflow
    platform.
\end{itemize}

A Canvasflow-managed mapping endpoint is a candidate feature for a future version when
tenant mapping volume or complexity grows beyond what can be maintained in an IdP hook.
